\section{A positive payment from the variation on a physical segment}
\label{sec:v110-variation}
The comparison pool of Section~\ref{sec:v109-pool} is left unchanged.
We first make a different use of the completed endpoint estimate: it
pays part of the \emph{original source}, whether or not that source
has a comparison window. All statements in this section are made
on the inherited one-copy physical pairing space.

Write $C_{\rm reg}$ for a fixed constant in
\eqref{eq:v109-regular-variation}. For fixed $\varepsilon>0$ define
\begin{equation}\label{eq:v110-variation}
 \begin{split}
 \mathcal V^\varepsilon(\theta)
   &=\operatorname{Var}_{[a(\theta),b(\theta)]}
                         f^\varepsilon(\,\cdot\,,\theta),\\
 \mathcal J(\theta)&=\int_{a(\theta)}^{b(\theta)}
               \frac{A(u,\theta)I_m(u,\theta)}{z(u,\theta)}\,du.
 \end{split}
\end{equation}
Endpoint values in the variation are the one-sided physical traces,
not values assigned on a continuation beyond the original event.

\begin{lemma}[Measurability and the full variation bound]
\label{lem:v110-variation}
Both functions in \eqref{eq:v110-variation} are Borel and finite on
every positive-length physical segment. Moreover,
\begin{equation}\label{eq:v110-full-variation}
 \mathcal V^\varepsilon
   =\int_a^b|\partial_u f^\varepsilon|\,du
   \le C_{\rm reg}(1+\varepsilon^{-1})\mathcal J.
\end{equation}
This is a bound for total variation, not merely for the difference
of the two endpoint values.
\end{lemma}
\begin{proof}
Pull the closed segment back to $[0,1]$ by $u=a+s(b-a)$, using the
endpoint traces at $s=0,1$. The resulting function is jointly Borel
and continuous in $s$. Its variation is the supremum of the sums of
absolute increments over finite rational partitions containing
$0$ and $1$. Continuity permits this countable family of partitions,
so its supremum is Borel. The integral defining $\mathcal J$ is
Borel by the same change of variables and nonnegative integration.
Theorem~\ref{thm:v109-traces} gives absolute continuity on the closed
segment. The variation of an absolutely continuous real function is
the integral of the absolute value of its derivative. Integrate
\eqref{eq:v109-regular-variation}; its right side is integrable by
Corollary~\ref{cor:v109-marked} and Lemma~\ref{lem:v109-gauss}.
This proves the equality, inequality and finiteness asserted above.
\end{proof}

For $\gamma\ge1$, split the original paired coefficient into
\begin{equation}\label{eq:v110-positive-split}
 c_{\rm var}^{\varepsilon,\gamma}
       =\min\{c^\varepsilon,\gamma\mathcal V^\varepsilon\},\qquad
 c_{\rm per}^{\varepsilon,\gamma}
       =(c^\varepsilon-\gamma\mathcal V^\varepsilon)_+.
\end{equation}
Thus $c^\varepsilon=c_{\rm var}^{\varepsilon,\gamma}
+c_{\rm per}^{\varepsilon,\gamma}$, with both terms nonnegative.
The subscript $\mathrm{per}$ will be justified by a physical
persistence statement in the next section. It denotes a positive
\emph{submeasure} of the original source, not a replacement of its
coefficient by a different normalization.

For a jointly Borel selector $0\le H\le1$ supported on $0<\ell<h$,
set
\begin{equation}\label{eq:v110-selected-sources}
 T_{H,j}^{\varepsilon,\gamma}(t)
   =\int c_j^{\varepsilon,\gamma}(\theta)
       \one_{\{a<t<b\}}H(\Psi(\theta,t))\,d\theta,
       \qquad j\in\{\mathrm{var},\mathrm{per}\}.
\end{equation}
The original selected paired source is their exact sum. The selector
is evaluated at $t$ and is not moved to another roof in the proof.

\begin{theorem}[A summed payment of the positive short source]
\label{thm:v110-payment}
At every finite collision count, on the original common roof versions,
\begin{equation}\label{eq:v110-finite-payment}
 0\le T_{H,\mathrm{var}}^{\varepsilon,\gamma}(t)
 \le C_{\rm reg}\gamma(1+\varepsilon^{-1})
                         \int_{t-h}^{t+h}v^0(u)\,du.
\end{equation}
Consequently, for every $\varepsilon,h>0$ and $\gamma\ge1$ fixed
before the collision limit,
\begin{equation}\label{eq:v110-ordered-payment}
 \mathcal H_M(T_{H,\mathrm{var}}^{\varepsilon,\gamma})
                  \le C\gamma(1+\varepsilon^{-1})h.
\end{equation}
The constants are uniform in the collision count, radius, exact
labels, segment length and selector. No positive lower bound for
$\ell$ or for an interior incidence is required.
\end{theorem}
\begin{proof}
For a segment meeting $t$ and satisfying $\ell<h$, one has
$(a,b)\subset(t-h,t+h)$. By \eqref{eq:v110-positive-split} and
\eqref{eq:v110-full-variation}, its selected contribution is at most
\[
 C_{\rm reg}\gamma(1+\varepsilon^{-1})
 \one_{\{a<t<b,\,\ell<h\}}H(\Psi(\theta,t))
             \int_a^b AI_m/z\,du.
\]
First sum these nonnegative quantities over the physical sheets.
Tonelli's theorem applies, without an estimate for the number of
sheets. At an integration roof $u$, each retained pair $(\theta,u)$
is still a point of the same original exact event. Dropping its
selector, which is between zero and one, and using physical
multiplicity one therefore bounds the sum by
\[
 C_{\rm reg}\gamma(1+\varepsilon^{-1})
 \int_{t-h}^{t+h}\int_{\{a(\theta)<u<b(\theta)\}}
              A(u,\theta)I_m(u,\theta)/z(u,\theta)\,d\theta\,du.
\]
The coarea identity in Input statement~\ref{input:v109-analytic}
identifies the inner integral with $v^0(u)$. This proves the
finite-count inequality. Apply \eqref{eq:v109-volume-average} at
the previously fixed width $h$, using
$\int_{t-h}^{t+h}v^0=2hA_hv^0$. Only after the complete sum and this
inequality are established are the uniform norms and collision
limsup taken. This proves \eqref{eq:v110-ordered-payment}.
\end{proof}

\begin{corollary}[The original microscopic schedule]
\label{cor:v110-schedule}
Keep $\widetilde H_B,h_B,r_B,K_B$ exactly as in
\eqref{eq:v109-original-micro} and \eqref{eq:v109-explicit-widths},
and choose
\begin{equation}\label{eq:v110-gamma}
 \gamma_B=\max\{2,\varepsilon(B)^{-1/32}\}.
\end{equation}
Let $\mathsf U_B^{\rm var}$ and $\mathsf U_B^{\rm per}$ denote the
two sources in \eqref{eq:v110-selected-sources} with
$H=\widetilde H_B$, $\varepsilon=\varepsilon(B)$ and
$\gamma=\gamma_B$. Then
\begin{equation}\label{eq:v110-source-identity}
 \mathsf U_B^{\rm micro}=\mathsf U_B^{\rm var}
                         +\mathsf U_B^{\rm per},\qquad
 \mathcal H_M(\mathsf U_B^{\rm var})\le C_MB^{-95/384}.
\end{equation}
\end{corollary}
\begin{proof}
The identity holds before pushforward, by
\eqref{eq:v110-positive-split}. For sufficiently large $B$,
$h_B=O(\varepsilon^4)$ and $\gamma_B=O(\varepsilon^{-1/32})$.
Theorem~\ref{thm:v110-payment} gives the bound
\[
 C\gamma_B(1+\varepsilon^{-1})h_B
      =O(\varepsilon^{95/32})=O(B^{-95/384}).
\]
All three quantities $h_B,\varepsilon(B),\gamma_B$ are fixed before
$m$ tends to infinity. No collision-dependent residence is inserted
into an estimate asserted only for fixed residence.
\end{proof}

Unlike \eqref{eq:v109-grazing-zero}, this payment does not rely on a
zero paired coefficient. If $c^\varepsilon>0$ and
$\mathcal V^\varepsilon>0$, its coefficient in
\eqref{eq:v110-positive-split} is positive. If
$c^\varepsilon\le\gamma\mathcal V^\varepsilon$, the complete
original coefficient is paid. These statements do not assert a
positive uniform proportion of paid mass in every exact event.
